\section{随体热模型与显热—潜热分配}
\label{sec:thermal}
本节全部使用物理单位。ALE 指材料坐标与空间坐标并用的描述：
冰的热状态在材料网格中演化，水的热状态在固定空间网格中演化，
两者通过当前位姿和覆盖率交换热量。
在不发生相变的区域，导热与对流的连续物理背景为
\begin{align}
 \rho_ic_i(\partial_t+\bm u_b\cdot\nabla)T_{\rm ice}&=\nabla\cdot(k_i\nabla T_{\rm ice}),\\
 \rho_wc_w(\partial_t+\bm u\cdot\nabla)T_{\rm water}&=\nabla\cdot(k_w\nabla T_{\rm water}).
\end{align}
随体坐标消去冰的刚体输运；相变所需热量通过后述局部能量分配扣除。
程序实际推进广延量和离散通量，界面附近使用其专门的覆盖率近似。

\subsection{广延状态与温度恢复}
每个冰材料单元保存固体质量 $m_a$ 和显热 $S_a$；
每个空间水单元保存面积 $V_i$ 和显热 $E_i$。
以熔点 $T_m$ 为显热零点，初值为
\begin{equation}
 m_a^0=\rho_iA,\quad S_a^0=m_a^0c_i(T_{i,0}-T_m),\quad
 V_i^0=A\phi_i^0\ \ (i\in\mathcal W),\quad
 E_i^0=\rho_wc_wV_i^0(T_{w,0}-T_m).
\end{equation}
温度由广延量恢复：
\begin{equation}
 T_a=\min\left(T_m,T_m+\frac{S_a}{m_ac_i}\right),\qquad
 T_i=T_m+\frac{E_i}{\rho_wc_wV_i}.
 \label{eq:temperature}
\end{equation}
仅在分母大于代码中的小量阈值时作除法。
空冰材料单元返回 $T_m$，无水单元返回配置中的初始空气温度。
使用质量与能量而非直接推进温度，可使融化、单元容量变化和能量收支共享同一组状态。

\subsection{冰内部导热}
对材料网格四邻接面，定义从 $a$ 到 $b$ 的导热功率
\begin{equation}
 \dot Q_{a\to b}=k_i\min(s_a,s_b)(T_a-T_b),\qquad
 \Delta S_a^{\rm cond}=-\delta t\sum_{b\sim a}\dot Q_{a\to b}.
\end{equation}
二维面长 $\Delta x$ 与中心距离 $\Delta x$ 相消，因此式中无额外网格因子。
两端在同一不可变质量与显热状态上计算，满足
$\dot Q_{a\to b}=-\dot Q_{b\to a}$。
\code{_body_heat_delta} 暂存导热增量，待加入冰水界面热量后统一更新材料。

\subsection{水内部导热与容器热边界}
以 $x$ 向面为例，左、右单元为 $L,R$，定义水占据比例
$a_i=\clip(V_i/A,0,1)$。两侧均为合法热水单元时
\begin{equation}
 F_{E,f}^{\rm cond}=-k_w\min(a_L,a_R)(T_R-T_L).
\end{equation}
约定面通量沿正坐标方向为正，则
\begin{equation}
 E_{ij}^{\rm new}=E_{ij}-\delta t
 (F^x_{i+1,j}-F^x_{i,j}+F^y_{i,j+1}-F^y_{i,j}).
 \label{eq:fv}
\end{equation}
同一面通量由两侧以相反符号使用，内部热交换在总和中抵消。
水—气面通量为零；冰水交换另行计算。

绝热容器壁的热通量为零。定温壁以半格距离 $\Delta x/2$ 计算入域热流
\begin{equation}
 \dot Q_{\partial\Omega\to i}=2k_wa_i(T_b-T_i).
\end{equation}
累计边界输入 $Q_\partial$ 由每个导热子步的入域功率乘以 $\delta t$ 累加。
定温边界只通过相邻水单元起作用，当前路径没有直接冰—容器导热项。
默认示例的四壁均为绝热。

\subsection{水体积与显热的成对对流}
面物理速度由两侧受力速度的算术平均给出：
\begin{equation}
 u_f=\frac{c_u}{2}(u_L^\ell+u_R^\ell).
\end{equation}
按 $u_f$ 的符号选上风单元 $D$，并同时构造体积与显热通量
\begin{equation}
 F_{V,f}=u_f\Delta x\,\clip(V_D/A,0,1),\qquad
 F_{E,f}^{\rm adv}=F_{V,f}\frac{E_D}{V_D}.
 \label{eq:paired-flux}
\end{equation}
两种广延量分别用式 \eqref{eq:fv} 的离散散度更新。
面两侧都必须属于当前湿流体集合，单元不能通过干燥气相或固体面直接输运水。
界面移动所需的容量调整由第 \ref{sec:remap} 节的重映射承担。
同温水满足 $E_i=eV_i$，成对通量在纯对流中维持这一比例，避免容量变化造成伪温变。

\src{simulator.py:_accumulate_body_conduction, _compute_water_conduction_flux_x, _compute_water_conduction_flux_y, _compute_water_advection_flux_x, _compute_water_advection_flux_y}

\subsection{冰水界面的热交换}
对每个湿流体单元 $i$，检查本单元及四个轴向邻单元的冰覆盖率。
每个被检查位置 $r$ 再映射至至多四个材料节点，故一个水线程最多形成 20 个热请求。
采用调和导热系数
\begin{equation}
 k_{wi}=\frac{2k_wk_i}{k_w+k_i},\qquad
 A_{ir}=\min\left(a_i,\min(1,C_r/s_*)\right).
\end{equation}
对仍有质量的材料节点 $a$，令双线性权重为 $W_{ra}$，请求热量为
\begin{equation}
 q_{ira}=\delta t\,k_{wi}A_{ir}\frac{W_{ra}s_a}{C_r}
 \max(0,T_i-T_a),\qquad C_r>0.
\end{equation}
这一接触系数是基于覆盖率的离散近似，未重建几何界面的精确面长。
为保证水只交出实际可用的显热，对同一水单元所有请求统一缩放：
\begin{equation}
 q_i=\sum_{r,a}q_{ira},\qquad
 \eta_i=\begin{cases}\min(1,\max(E_i,0)/q_i),&q_i>0,\\0,&q_i=0,\end{cases}
 \qquad \widehat q_{ira}=\eta_iq_{ira}.
\end{equation}
水显热减少 $\sum_{r,a}\widehat q_{ira}$，相应材料热增量增加同量。
材料端采用原子累加，水线程只更新自身显热。
该热交换只允许水向冰传热；与单向融化假设一致。
当冰与热水尚无离散接触时，空气不会为冰供热。

\subsection{局部相变更新}
设材料单元累计得到净热量 $Q_a$，其旧状态为 $(m_a,S_a)$，潜热为 $L$。
若 $Q_a\leq0$ 且仍有固体，则 $S_a' = S_a+Q_a$、$m_a'=m_a$。
若 $Q_a>0$，先补足低于熔点的显热亏损，再支付潜热：
\begin{align}
 Q_{\rm sens}&=\min[Q_a,\max(-S_a,0)],& S_a'&=S_a+Q_{\rm sens},\\
 Q_{\rm rest}&=Q_a-Q_{\rm sens},&
 \Delta m_a&=\min(m_a,Q_{\rm rest}/L),\\
 m_a'&=m_a-\Delta m_a,&
 E_a^{\rm melt}&=\max(0,Q_{\rm rest}-L\Delta m_a).
 \label{eq:melt-rule}
\end{align}
材料完全融化后，其 $m_a',S_a'$ 置零，剩余正热量 $E_a^{\rm melt}$ 随融水进入空间网格。
因此，即使剩余质量极小，也必须由实际潜热支付才能消失。
这是显热与潜热分配实现的焓法；并不存在另外一个被时间积分的统一温度—液相率场。

\subsection{融水注入与密度转换}
\begin{equation}
 \Delta V_w=\Delta m_a/\rho_w,
 \qquad \Delta E_w=E_a^{\rm melt},\qquad
 \frac{\Delta V_w}{\Delta V_i}=\frac{\rho_i}{\rho_w}.
 \label{eq:melt-volume}
\end{equation}
潜热记入累计融化能，不重复加入水显热。
材料节点优先向已记录的界面水节点注入；若有多个供热节点，
记录其线性编码的最大值，以原子最大值操作确定一个目标。
无记录目标时，在材料位置附近 $5\times5$ 网格中寻找最近合法湿单元。
资格掩码先由单独内核冻结，防止注入过程改变后续线程的搜索集合。

仍无法分配的融水进入后备池，按合法湿单元的剩余容量分配；
无剩余容量权重时按已有水体积分配，再由容量同步恢复单元限制。
完全没有合法湿单元则报错。
后备分配保持质量与显热总量，但不代表精确的局部融水轨迹。

\src{simulator.py:_exchange_interface_heat, _apply_body_heat_and_phase_change, _freeze_melt_injection_eligibility, _inject_melt_water, _distribute_unassigned_melt}

\subsection{快慢时间尺度与稳定性}
\label{sec:multirate}
每个 LBM 步执行一次快速过程：容量重映射及水的对流。
每累计 $K$ 个 LBM 步执行一次慢过程：导热、边界热交换、融化和融水注入。
默认 $K=8$；慢过程按当前固定刚体位姿推进 $K\Delta t$，不是沿区间内全部历史位姿积分。

对流步长依据\textbf{面向外速度之和}而非单个节点速度范数。
对合法湿面，令格子面速度为 $\bm u_f^\ell$，外法向为 $\bm n_{if}$，定义
\begin{equation}
 r_i=\sum_{f\subset\partial i}\max(0,\bm u_f^\ell\cdot\bm n_{if}),
 \quad r_{\max}=\max_i r_i,
 \quad \delta t_{\rm adv}^{\max}=\frac{C_{\max}\Delta t}{r_{\max}}.
\end{equation}
无流出或关闭对流时上界为无穷。导热步长上界为
\begin{equation}
 \alpha_{\max}=\max\left(\frac{k_w}{\rho_wc_w},\frac{k_i}{\rho_ic_i}\right),
 \qquad \delta t_{\rm diff}^{\max}=Fo_{\max}\frac{\Delta x^2}{\alpha_{\max}}.
\end{equation}
给定过程区间 $\Delta T$，取
\begin{equation}
 N_{\rm sub}=\max\left(1,\left\lceil\Delta T/\delta t^{\max}-10^{-14}\right\rceil\right),
 \qquad \delta t=\Delta T/N_{\rm sub}.
\end{equation}
默认 $C_{\max}=0.5$、$Fo_{\max}=0.15$，每次过程至多 64 个子步。
配置要求 $Fo_{\max}\leq1/6$，与定温壁半格距离和角点扩散系数一致。
上述是代码实际采用的显式步长控制，不能替代整体耦合方法的收敛性验证。
慢热时钟只在慢过程结束时增加；若直接调用非 $K$ 倍数的 \code{step}，
其热时钟可落后于流动时钟。
